% ============================================================================
% Dok. 265 -- Korrespondenzebenen zwischen physikalischen Formalismen:
%             Blume-Maleev / FFGFT als Arbeitsbeispiel
%
% Zweck: Arbeitsbeispiel für den IPI-Rahmen zum strukturellen Modellvergleich
%        (Stefaan Vossen, José Guevara RA-Architektur, Juni 2026).
%        Zeigt drei unterschiedliche Korrespondenzebenen zwischen zwei
%        unabhängig hergeleiteten Formalismen am konkreten Testfall.
%
% Status: Methodisches Dokument. Keine neue FFGFT-Physik. Die FFGFT-
%         Interpretationen stammen aus Dok. 264; dieses Dokument fügt die
%         Metaklassifikation hinzu, welchen Korrespondenztyp jeder Punkt
%         darstellt.
% ============================================================================

\section*{Zweck}

Dieses Dokument liefert ein Arbeitsbeispiel für den strukturellen Vergleich
physikalischer Formalismen, entwickelt für die Modell-Einordnungsarbeit der
IPI-Gruppe. Das Beispiel verwendet zwei unabhängig hergeleitete mathematische
Beschreibungen desselben experimentellen Systems -- den ferrotoroidischen
Quaternärspeicher LiNi$_{0.8}$Fe$_{0.2}$PO$_4$ (Qureshi et al.\ 2026) --
um zu illustrieren, dass ,,Korrespondenz zwischen Formalismen`` keine einzelne
Kategorie ist, sondern eine geschichtete Relation mit mindestens drei
verschiedenen Ebenen.

Die zwei Beschreibungen sind:
\begin{itemize}
  \item \textbf{Blume-Maleev-Polarisationsmatrix-Formalismus}: die Standard-
        Festkörperbeschreibung, vollständig und quantitativ korrekt, vollständig
        innerhalb etablierter experimenteller Physik entwickelt.
  \item \textbf{FFGFT (T$^4$-Torus-Geometrie)}: das geometrische Feldtheorie-
        Rahmenwerk, auf dasselbe System aus eigenen ersten Prinzipien angewendet,
        ohne Anleihen bei der Festkörperbeschreibung.
\end{itemize}

\noindent Keines der Rahmenwerke wurde im Hinblick auf das andere entwickelt.
Die auftretenden Korrespondenzen sind daher strukturell informativ -- sie
zeigen, welcher mathematische Inhalt geteilt wird und welcher
rahmenwerks-spezifisch ist.

% ============================================================================
\section{Die drei Korrespondenzebenen}
% ============================================================================

\begin{table}[htbp]
\centering
\caption{Klassifikation der Korrespondenztypen zwischen Blume-Maleev-
und FFGFT-Beschreibung des ferrotoroidischen Systems.}
\label{tab:ebenen}
\begin{tabular}{p{2.5cm}p{5.5cm}p{5.5cm}}
\toprule
\textbf{Ebene} & \textbf{Definition} & \textbf{Diagnosekriterium} \\
\midrule
\textbf{1. Fast-Äquivalenz} &
  Beide Formalismen drücken dasselbe mathematische Objekt aus, unabhängig
  aus verschiedenen Ausgangspunkten hergeleitet. &
  Die abstrakte Struktur (Gruppe, Operator, Invariante) ist identisch;
  nur Herleitungsweg und physikalische Interpretation unterscheiden sich. \\
\addlinespace
\textbf{2. Struktureller Überlapp} &
  Dieselbe qualitative Struktur erscheint in beiden, aber der quantitative
  Inhalt wird nicht geteilt. &
  Ein Formalismus kann für das strukturelle Gerüst auf den anderen
  abgebildet werden, nicht aber für numerische Vorhersagen. \\
\addlinespace
\textbf{3. Komplementär} &
  Jeder Formalismus erklärt etwas, das der andere nicht kann. &
  Keiner ist Teilmenge oder Erweiterung des anderen; die Vereinigung
  deckt mehr ab als jeder allein. \\
\bottomrule
\end{tabular}
\end{table}

% ============================================================================
\section{Der Vergleichsoperator -- explizit gemacht}
% ============================================================================

Die Drei-Ebenen-Klassifikation in Tabelle~\ref{tab:ebenen} weist jede
Korrespondenz einer Ebene zu. Bis hierher blieb jedoch \emph{implizit},
\emph{nach welchen Kriterien} diese Zuordnung erfolgt -- welcher Operator
entscheidet, ob ein Phänomen auf Ebene~1, 2 oder~3 fällt. Dieser Abschnitt
legt diesen Operator offen, sodass er prüfbar und bestreitbar wird statt
auf Intuition zu beruhen.

\subsection{Die sechs Vergleichsachsen}

Zwei Formalismen werden entlang von sechs unabhängigen Achsen verglichen.
Auf jeder Achse ist der Inhalt entweder \emph{geteilt} (beide Formalismen
liefern dasselbe) oder \emph{nicht geteilt}:

\begin{enumerate}
  \item \textbf{Geteilte Invarianten}: Erzeugen beide dieselben
        Erhaltungsgrößen / topologischen Invarianten?
        (z.\,B.\ die erhaltene Windungs-/Domänenzahl)
  \item \textbf{Geteilte Operatoren}: Tritt dieselbe abstrakte
        Operatorstruktur auf? (z.\,B.\ dieselbe Gruppe
        $\mathbb{Z}_2\times\mathbb{Z}_2$, dieselbe Transformationsregel)
  \item \textbf{Geteilte mathematische Strukturen}: Teilen beide das
        qualitative Gerüst -- antisymmetrische Kopplung, Projektionsfaktor --
        auch wenn die Zahlenwerte abweichen?
  \item \textbf{Geteilte Vorhersagen}: Liefern beide für dasselbe
        Experiment denselben \emph{quantitativen} Wert?
  \item \textbf{Geteilte physikalische Mechanismen}: Schreiben beide das
        Phänomen \emph{derselben} physikalischen Ursache zu?
  \item \textbf{Geteilte Hardware-Implikationen}: Folgen aus beiden
        dieselben Konsequenzen für reale Bauelemente / Messungen?
\end{enumerate}

\subsection{Wie die Achsen die Ebenen definieren}

Die Ebenenzuordnung ist keine Intuition, sondern liest sich direkt aus
dem Muster der geteilten Achsen ab:

\begin{table}[htbp]
\centering
\caption{Der Vergleichsoperator: welche Achsen auf welcher Ebene geteilt
sind. $\checkmark$ = geteilt, $\times$ = nicht geteilt, $\sim$ = nur
qualitativ geteilt.}
\label{tab:operator}
\begin{tabular}{p{5.5cm}ccc}
\toprule
\textbf{Vergleichsachse} & \textbf{Ebene 1} & \textbf{Ebene 2} & \textbf{Ebene 3} \\
 & Fast-Äquiv. & Überlapp & Komplement. \\
\midrule
1. Invarianten          & $\checkmark$ & $\sim$       & $\times$ \\
2. Operatoren           & $\checkmark$ & $\sim$       & $\times$ \\
3. Mathematische Struktur & $\checkmark$ & $\checkmark$ & $\times$ \\
4. Vorhersagen (quantitativ) & $\times$  & $\times$     & $\times$ \\
5. Physikalische Mechanismen & $\times$  & $\times$     & $\times$ \\
6. Hardware-Implikationen & $\times$    & $\times$     & $\times$ \\
\bottomrule
\end{tabular}
\end{table}

\noindent Das Zuordnungskriterium lautet damit explizit:
\begin{itemize}
  \item \textbf{Ebene 1 (Fast-Äquivalenz):} Achsen 1--3 voll geteilt --
        identische Invariante, identischer Operator, identische Struktur --
        bei verschiedenem Herleitungsweg. Achsen 4--6 bleiben getrennt.
        (Beispiele: $\mathbb{Z}_2\times\mathbb{Z}_2$, $\mathcal{P}''=0$,
        topologische Ladung.)
  \item \textbf{Ebene 2 (Struktureller Überlapp):} Achse~3 geteilt
        (dasselbe qualitative Gerüst), Achsen 1--2 nur teilweise, Achsen
        4--6 nicht. (Beispiele: antisymmetrischer Austausch,
        Dekompaktifizierungsprojektion.)
  \item \textbf{Ebene 3 (Komplementär):} Keine der Achsen 1--3 geteilt;
        jeder Formalismus deckt einen Bereich ab, den der andere nicht
        erreicht. (Beispiele: geometrische Notwendigkeit von
        $\mathbb{Z}_2\times\mathbb{Z}_2$ in FFGFT; materialspezifische
        Übergangstemperaturen in Blume-Maleev.)
\end{itemize}

Auffällig: Achsen~4--6 (quantitative Vorhersagen, physikalische
Mechanismen, Hardware) sind in diesem Vergleich \emph{auf keiner} Ebene
geteilt. Das ist kein Mangel der Klassifikation, sondern ihr Hauptbefund:
Blume-Maleev und FFGFT teilen mathematische \emph{Struktur} (Achsen 1--3),
aber nicht physikalische \emph{Substanz} (Achsen 4--6). Genau das
unterscheidet diese Korrespondenz von einer echten Äquivalenz
(siehe Abschnitt zur Heisenberg/Schrödinger-Analogie).

\subsection{Residuen: Grenzen des Operators}

Der so explizit gemachte Operator ist selbst nicht vollständig
ausgereift. Drei Residuen bleiben offen und werden hier deklariert,
nicht verborgen:

\begin{itemize}
  \item \textbf{Achsen-Vollständigkeit:} Sind sechs Achsen ausreichend,
        oder fehlen Dimensionen (z.\,B.\ geteilte Randbedingungen,
        geteilte Symmetriebrechungsmuster)? Die Liste ist ein Arbeitsstand,
        keine bewiesene Basis.
  \item \textbf{Achsen-Gewichtung:} Alle sechs Achsen werden hier
        gleichgewichtet. Ob eine geteilte Invariante (Achse~1) schwerer
        wiegen sollte als eine geteilte Struktur (Achse~3), ist nicht
        entschieden.
  \item \textbf{Achsen-Unabhängigkeit:} Die Achsen sind als unabhängig
        behandelt, sind es aber möglicherweise nicht -- eine geteilte
        Operatorstruktur (Achse~2) impliziert oft eine geteilte Invariante
        (Achse~1). Die Korrelationsstruktur der Achsen ist nicht analysiert.
\end{itemize}

Diese Residuen ändern die Ebenenzuordnung im vorliegenden Fall nicht
(das Muster in Tabelle~\ref{tab:operator} ist robust), markieren aber,
wo der Operator selbst weiter zu prüfen ist. Der Wert der Explizitmachung
liegt gerade darin: ein zuvor implizites, aber tragendes Kriterium wird
zu einem prüfbaren Objekt -- inspizierbar statt überzeugend.

% ============================================================================
\section{Ebene 1 -- Fast-Äquivalenzen}
% ============================================================================

\subsection{Die $\mathbb{Z}_2\times\mathbb{Z}_2$-Domänenstruktur}

\textbf{In Blume-Maleev:} Die vier magnetischen Domänen sind durch zwei
Symmetrieoperationen des Kristalls verknüpft: Zeitumkehr $1'$ (bildet
$a \leftrightarrow b$ ab, der ferrotoroidische Flip) und die zweizählige
Schraubenachse $2_1\|y$ (bildet $'\leftrightarrow''$ ab, der Orientierungsflip).
Zusammen erzeugen diese eine $\mathbb{Z}_2\times\mathbb{Z}_2$-Gruppenstruktur
über den vier Domänen, hergeleitet aus der Kristallsymmetriegruppe P$1121/a'$.

\textbf{In FFGFT:} Der T$^4$-Torus hat zwei unabhängige binäre Freiheitsgrade
für geschlossene Windungsschleifen: die zeitliche Windungsrichtung ($\pm\tilde{T}$,
ob die Schleife im oder gegen den Uhrzeigersinn in der kompakten Zeitrichtung
windet) und die räumliche Windungsparität ($\pm c$, die zwei möglichen
Orientierungen der Windungsebene im beobachtbaren 3D-Unterraum). Diese
erzeugen dieselbe $\mathbb{Z}_2\times\mathbb{Z}_2$-Gruppenstruktur über den
vier Windungsklassen.

\textbf{Bewertung:} Dieselbe abstrakte Gruppe $\mathbb{Z}_2\times\mathbb{Z}_2$
wird aus zwei völlig verschiedenen Ausgangspunkten hergeleitet -- Kristall-
symmetrie auf der einen Seite, T$^4$-Kompaktifizierungsgeometrie auf der
anderen. Das ist eine Fast-Äquivalenz: das mathematische Objekt ist dasselbe,
die Herleitungswege sind unabhängig.

\textbf{Was das nicht bedeutet:} Das impliziert nicht, dass das Kristallgitter
\emph{ein} T$^4$-Torus \emph{ist}, oder dass die FFGFT-Parameter den
spezifischen Kristall beschreiben. Die geteilte Gruppenstruktur ist eine
mathematische Tatsache darüber, wie sich binäre Symmetrien kombinieren --
sie identifiziert nicht den zugrundeliegenden physikalischen Mechanismus.

\subsection{Die $\mathcal{P}'' = 0$-Bedingung}

\textbf{In Blume-Maleev:} Die allgemeine Polarisationstransformation ist
$\mathbf{P}_f = \mathcal{P}'\mathbf{P}_i + \mathcal{P}''$. Für
LiNi$_{0.8}$Fe$_{0.2}$PO$_4$ gilt $\mathcal{P}'' = 0$, weil die magnetischen
Strukturformfaktoren rein imaginär sind: $\mathrm{Re}(NM_{\perp i}^*) = 0$
und $\mathrm{Im}(M_{\perp y}M_{\perp z}^*) = 0$. Dies folgt aus der
antiferromagnetischen Head-to-Tail-Konfiguration.

\textbf{In FFGFT:} Der Rekursionsoperator $\mathcal{R}$ (Dok.~203) auf dem
T$^4$-Modenraum zerfällt in einen rotatorischen Anteil (kohärente
Modentransformation, Windungszahl erhalten) und einen Erzeugungsanteil
(spontane Modengenerierung an T$^4$-Grenzknoten, Dok.~204). Im statischen
T$^4$-Grundzustand -- keine Expansion, keine Horizont-Dynamik -- verschwindet
der Erzeugungsanteil: $\mathcal{P}'' = 0$.

\textbf{Bewertung:} Beide leiten dasselbe Verschwinden des Erzeugungs-/
Vernichtungsterms aus einer ,,keine spontane Erzeugung``-Bedingung ab --
im Kristall aus der antiferromagnetischen Symmetrie, in FFGFT aus der
statischen Geometrie. Die mathematische Struktur ist dieselbe.

\subsection{Erhaltung der topologischen Ladung}

\textbf{Im Paper:} Der Speicher ist nichtflüchtig -- Domänenpopulationen
bleiben bei $E = H = 0$ nach der Feldkühlung stabil. Das Paper identifiziert
dies als Konsequenz der Energiebarriere zwischen toroidalen Domänen.

\textbf{In FFGFT:} Ein Modus mit von Null verschiedener Windungszahl $n$
kann nicht auf $n = 0$ zerfallen, ohne die Mindestenergie der T$^4$-Rand-
bedingung zu überwinden (Dok.~204). Die Windungszahl ist eine topologische
Invariante.

\textbf{Bewertung:} Beide beschreiben dasselbe mathematische Phänomen --
eine diskrete topologische Invariante, geschützt durch eine Energiebarriere.
Fast-Äquivalenz auf der Ebene des Erhaltungsgesetzes; verschiedene Mechanismen
auf der Ebene des physikalischen Ursprungs.

% ============================================================================
\section{Ebene 2 -- Struktureller Überlapp ohne Äquivalenz}
% ============================================================================

\subsection{Antisymmetrischer Austausch und das Nichtabschluss-Theorem}

\textbf{In Blume-Maleev:} Die Dzyaloshinskii-Moriya-Wechselwirkung ist
$\mathcal{H}^{\mathrm{DM}} = \mathbf{D}\cdot(\mathbf{S}_1\times\mathbf{S}_2)$,
ungerade unter Spinaustausch (antisymmetrisch), mit $\mathbf{D}_{34} = -\mathbf{D}_{12}$
aus der Inversionssymmetrie.

\textbf{In FFGFT:} Das Nichtabschluss-Theorem (Dok.~193) besagt, dass
Kreuzprodukte von Modenlabeln aus \emph{verschiedenen} Windungsklassen
topologisch unzulässig sind. Das antisymmetrische Kreuzprodukt $S_i\times S_j$
entspricht strukturell diesem Befund: es koppelt zwei Spinkomponenten (zwei
verschiedene Windungsrichtungen) in einer Weise, die antisymmetrisch unter
Vertauschung ist.

\textbf{Bewertung:} Die antisymmetrische Kreuzproduktstruktur wird geteilt.
Aber die DM-Vektorkomponenten $D^b$ und $D^c$ sind materialspezifische
Größen ohne FFGFT-Entsprechung. FFGFT kann sagen, dass die antisymmetrische
Struktur topologisch eingeschränkt ist; es kann die spezifische DM-Vektor-
orientierung oder -größe für diesen Kristall nicht vorhersagen.

\subsection{Dekompaktifizierungsprojektion}

\textbf{In Blume-Maleev:} Die Faktoren $p_{fx}/p_{ix}$ messen, wie genau
der Messapparat den Polarisationszustand auflöst.

\textbf{In FFGFT:} Die Dekompaktifizierungsprojektion (Schicht~1 zu Schicht~2,
Dok.~001) misst den Anteil des einheitenlosen geometrischen Inhalts, der in
der SI-projizierten Beschreibung beobachtbar wird.

\textbf{Bewertung:} Die Struktur -- ein Projektionsfaktor, der zwischen dem
abstrakten mathematischen Inhalt und der beobachtbaren Größe vermittelt -- ist
dieselbe in beiden. Aber die spezifischen Projektionsfaktoren $p_{fi}$ sind
apparatur- und materialabhängige Zahlen; in FFGFT wird die Projektion durch
$\xi = 4/30000$ und die SI-Einbettung bestimmt. Keine quantitative Abbildung.

% ============================================================================
\section{Ebene 3 -- Komplementär (Nicht-Überlappend)}
% ============================================================================

\subsection{Was FFGFT erklärt, was Blume-Maleev nicht erklärt}

Der Blume-Maleev-Formalismus leitet die Vier-Domänen-Struktur aus der
beobachteten Kristallsymmetrie ab. Er nimmt die Symmetriegruppe als gegeben
und arbeitet die Konsequenzen aus. Er erklärt nicht, \emph{warum} die
relevante Symmetriegruppe $\mathbb{Z}_2\times\mathbb{Z}_2$ ist und nicht eine
andere -- das ist durch das spezifische Material bestimmt und wird dem
Experiment entnommen.

FFGFT liefert eine geometrische Erklärung: $\mathbb{Z}_2\times\mathbb{Z}_2$
ist die \emph{notwendige} Struktur für einen kompakten T$^4$-Torus mit je
einer zeitartigen und einer raumartigen kompakten Richtung, die zur
beobachtbaren Domänenklassifikation beitragen. Die Vier-Zustands-Struktur
ist nicht kontingent auf das Material; sie ist eine Konsequenz der Topologie.

\subsection{Was Blume-Maleev liefert, was FFGFT nicht liefert}

Das Paper liefert, was FFGFT nicht kann: quantitative Vorhersagen für den
spezifischen Kristall. Die Übergangstemperaturen $T_2 = 25$\,K und
$T_4 = 21$\,K, die Domänenpopulationsgrade (bis zu 90\% Mehrheitsdomäne bei
optimaler Feldkühlung), die DM-Vektorausrichtung ($35°$ zur $c$-Achse), die
Polarisationstensorelemente in Tabelle~1 -- all das folgt aus den
Materialparametern von LiNi$_{0.8}$Fe$_{0.2}$PO$_4$ und liegt außerhalb des
FFGFT-Anwendungsbereichs.

% ============================================================================
\section{Die Analogie: Warum nicht Heisenberg vs.\ Schrödinger?}
% ============================================================================

Eine naheliegende erste Analogie für zwei Formalismen, die dieselbe Physik
beschreiben, ist Heisenbergs Matrizenmechanik vs.\ Schrödingers
Wellenmechanik (1925--26), die nachträglich als vollständig äquivalent
bewiesen wurden (von Neumann 1932). Aber diese Analogie ist für den
vorliegenden Fall zu stark.

Heisenberg und Schrödinger beschrieben \emph{dieselben physikalischen Objekte}
(Atome, Elektronen) auf \emph{derselben Skala} mit \emph{demselben Vorhersage-
umfang}. Ihre Äquivalenz bedeutete, dass jede Vorhersage des einen auch eine
Vorhersage des anderen war.

Die Blume-Maleev/FFGFT-Korrespondenz ist schwächer: die beiden Formalismen
teilen mathematische Struktur auf Ebenen~1 und~2, operieren aber auf
verschiedenen Skalen (Angström-Gitter vs.\ Sub-Planck-T$^4$), mit
verschiedenem Vorhersageumfang (materialspezifisch-quantitativ vs.\
topologisch-strukturell) und verschiedenen Herleitungsgrundlagen.

Eine präzisere Analogie ist die Relation zwischen:
\begin{itemize}
  \item Einer \textbf{Symmetriegruppen-Klassifikation} (z.\,B.\ Wigners
        Klassifikation von Teilchen durch irreduzible Darstellungen der
        Poincaré-Gruppe): liefert die Topologie, die Quantenzahlen, die
        Erhaltungsgesetze.
  \item Einer \textbf{Lagrange-Feldtheorie für ein spezifisches Teilchen}
        (z.\,B.\ QED für das Elektron): liefert alle quantitativen Vorhersagen
        für dieses spezifische Teilchen.
\end{itemize}

Die Gruppenklassifikation ist keine Teilmenge von QED, und QED ist keine
Teilmenge der Wigner-Klassifikation. Sie sind komplementär: die Gruppe liefert
die strukturelle Notwendigkeit (warum das Elektron existiert und Spin 1/2 hat);
die Feldtheorie liefert die Dynamik (wie es wechselwirkt).

Blume-Maleev entspricht der Feldtheorie (quantitativ, materialspezifisch,
dynamisch vollständig). FFGFT entspricht der Gruppenklassifikation (strukturelle
Notwendigkeit, skalenunabhängig, topologisch erklärend). Der geteilte Inhalt
auf Ebenen~1 und~2 ist der Schnitt -- die topologische Struktur, die beide
Beschreibungen enthalten müssen, weil sie eine Eigenschaft des physikalischen
Systems ist, unabhängig vom verwendeten Formalismus.

% ============================================================================
\section{Folgerungen für den Modellvergleich}
% ============================================================================

Die Drei-Ebenen-Gliederung hat praktische Konsequenzen für den Vergleich von
Rahmenwerken in der IPI-Gruppe:

\begin{enumerate}
  \item \textbf{Ebene-1-Korrespondenzen sind am informativsten}: Wenn zwei
        unabhängig hergeleitete Formalismen dasselbe mathematische Objekt
        erzeugen (dieselbe Gruppe, dieselbe Invariante, dasselbe Erhaltungsgesetz),
        ist das Evidenz dafür, dass das Objekt ein echter struktureller
        Charakterzug des physikalischen Systems ist -- kein Artefakt des
        Formalismus.
  \item \textbf{Ebene-2-Korrespondenzen definieren das gemeinsame Vokabular}:
        Der strukturelle Überlapp zeigt, welche Konzepte zwischen Rahmenwerken
        übertragbar sind (antisymmetrische Kopplung, Projektionsfaktoren) und
        welche nicht.
  \item \textbf{Ebene-3 definiert die echte Grenze}: Der exklusive Inhalt
        jedes Rahmenwerks markiert, wo das andere es nicht ersetzen kann.
        Das verhindert sowohl Übertreibung (,,mein Rahmenwerk erklärt alles,
        was das andere tut``) als auch Untertreibung (,,die Rahmenwerke haben
        nichts miteinander zu tun``).
  \item \textbf{Volle Äquivalenz ist der Spezialfall, nicht der Standard}:
        Die Heisenberg/Schrödinger-Äquivalenz war ein bemerkenswertes Ergebnis,
        keine Vorlage. Die meisten Rahmenwerke, die strukturellen Inhalt teilen,
        sind durch partiellen Überlapp oder Komplementarität verknüpft,
        nicht durch Äquivalenz.
\end{enumerate}

\noindent\textbf{Referenz.} Das als Testfall verwendete Ferrotoroidizitäts-
Experiment: N.~Qureshi et al., \textit{Toroidizität als Weg zu nichtflüchtigem
Quaternärspeicher in Antiferromagneten}, \textit{Nature Communications}
\textbf{17}, 4033 (2026). DOI:~10.1038/s41467-026-70767-8. Die FFGFT-
Interpretation aller Formeln des Papers befindet sich in Dok.~264 dieser Reihe.
